% ============================================================
% Paper 93 (DE): Gromovs Füllungsradius und Systolische Geometrie
% Autor: Michael Fuhrmann
% Projekt: Specialist Mathematics Project
% Build: 169
% Datum: 2026-03-12
% ============================================================
\documentclass[12pt,a4paper]{amsart}

\usepackage{amsmath,amssymb,amsthm,mathtools}
\usepackage{enumitem}
\usepackage{booktabs}
\usepackage{array}
\usepackage[hidelinks]{hyperref}
\usepackage[utf8]{inputenc}
\usepackage[T1]{fontenc}
\usepackage[ngerman]{babel}
\usepackage{geometry}
\geometry{margin=2.5cm}

% ---- Theorem-Umgebungen -------------------------------------
\newtheorem{theorem}{Theorem}[section]
\newtheorem{lemma}[theorem]{Lemma}
\newtheorem{corollary}[theorem]{Korollar}
\newtheorem{proposition}[theorem]{Proposition}
\theoremstyle{definition}
\newtheorem{definition}[theorem]{Definition}
\newtheorem{example}[theorem]{Beispiel}
\theoremstyle{remark}
\newtheorem{remark}[theorem]{Bemerkung}
\newtheorem{conjecture}[theorem]{Vermutung}
% Deutsche Aliasse
\newtheorem{satz}[theorem]{Satz}
\newtheorem{korollar}[theorem]{Korollar}
\newtheorem{vermutung}[theorem]{Vermutung}
\newtheorem{bemerkung}[theorem]{Bemerkung}
\newtheorem{definition_de}[theorem]{Definition}

% ---- Operatoren ---------------------------------------------
\DeclareMathOperator{\sys}{sys}
\DeclareMathOperator{\FillRad}{FüllRad}
\DeclareMathOperator{\vol}{vol}
\DeclareMathOperator{\diam}{diam}
\newcommand{\RR}{\mathbb{R}}
\newcommand{\ZZ}{\mathbb{Z}}
\newcommand{\RP}{\mathbb{RP}}
\newcommand{\TT}{\mathbb{T}}

% ---- Zusammenfassung ----------------------------------------
\begin{abstract}
Die systolische Geometrie untersucht die Beziehung zwischen der kürzesten
nicht-kontrahierbaren Schleife (der Systole) und dem Volumen einer
Riemannschen Mannigfaltigkeit.
Gromovs grundlegende Ungleichung $\sys(M)^n \le C_n \vol(M)$ gilt für alle
\emph{wesentlichen} $n$-Mannigfaltigkeiten.
Die entscheidende Zwischengröße ist der \emph{Füllungsradius} $\FillRad(M)$,
für den Gromov $\FillRad(M) \le C_n \sys(M)$ bewiesen hat.
Wir behandeln die klassischen Ungleichungen von Pu (für $\RP^2$) und Loewner
(für $\TT^2$), Gromovs Kompaktheits\-satz für metrische Räume sowie die
Rolle der Morse-Theorie für das Längenfunktional.
Scharfe systolische Konstanten sind nur in Spezialfällen bekannt;
viele bleiben offen.
\end{abstract}

\title{Gromovs Füllungsradius und Systolische Geometrie}

\author{Michael Fuhrmann}
\address{Specialist Mathematics Project, Build 169}
\email{specialist-maths@project.local}
\date{2026-03-12}

\maketitle
\tableofcontents

% =============================================================
\section{Einleitung}
% =============================================================

Eine der ältesten Fragen der Riemannschen Geometrie lautet: Wie kurz kann
eine geschlossene Kurve durch topologische Zwänge erzwungen werden?
Gegeben eine geschlossene Riemannsche Mannigfaltigkeit $(M,g)$, ist die
\emph{Systole} $\sys(M,g)$ die kürzeste Länge unter allen nicht-kontrahierbaren
geschlossenen Kurven in $M$.

Carl Loewner bewies 1949 (unveröffentlicht, zitiert in \cite{Pu1952}),
dass für jede Metrik $g$ auf dem 2-Torus $\TT^2$:
\[
  \sys(\TT^2, g)^2 \le \frac{2}{\sqrt{3}} \vol(\TT^2, g),
\]
mit Gleichheit für den flachen hexagonalen Torus.  Pu \cite{Pu1952} bewies
eine analoge Ungleichung für $\RP^2$.

1983 bewies Gromov \cite{Gromov1983} eine weitreichende Verallgemeinerung:
Für jede \emph{wesentliche} geschlossene $n$-Mannigfaltigkeit $M$ gilt
\[
  \sys(M)^n \le C_n \vol(M).
\]
Der Beweis führte eine neue geometrische Größe ein: den \emph{Füllungsradius}
$\FillRad(M)$.

% =============================================================
\section{Die Systole}
% =============================================================

\begin{definition_de}[Systole]
Sei $(M,g)$ eine geschlossene Riemannsche Mannigfaltigkeit mit
$\pi_1(M) \neq 1$.  Die \emph{1-Systole} ist
\[
  \sys_1(M,g) = \inf\bigl\{ \mathrm{Länge}(\gamma) :
    \gamma \text{ ist eine geschlossene nicht-kontrahierbare Schleife in } M
  \bigr\}.
\]
Das Infimum wird von einer Geodäten realisiert, wenn $M$ kompakt ist.
\end{definition_de}

\begin{example}
\begin{enumerate}[label=(\roman*)]
  \item Für den flachen Torus $\TT^2 = \RR^2/\Lambda$:
        $\sys = \lambda_1(\Lambda)$ (kürzester Gittervektor).
  \item Für $\RP^n$ mit der runden Metrik vom Kreisradius $R$:
        $\sys = \pi R$ (die projektive Gerade).
  \item Für eine hyperbolische Fläche vom Geschlecht $g \ge 2$:
        $\sys \le 2\log(4g-2)$.
\end{enumerate}
\end{example}

% =============================================================
\section{Der Füllungsradius}
% =============================================================

\begin{definition_de}[Füllungsradius \cite{Gromov1983}]
Sei $M$ eine geschlossene orientierbare $n$-Mannigfaltigkeit, isometrisch
eingebettet in den Banach-Raum $L^\infty(M)$ via der Kuratowski-Einbettung
$\iota \colon x \mapsto d(x,\cdot)$.  Der \emph{Füllungsradius} ist
\[
  \FillRad(M) = \inf\bigl\{ r > 0 : \iota_*[M] = 0 \in H_n(U_r(M);\ZZ)
  \bigr\},
\]
wobei $U_r(M) = \{f \in L^\infty(M) : d(f, \iota(M)) < r\}$ die
$r$-röhrenförmige Umgebung von $\iota(M)$ in $L^\infty(M)$ ist.
\end{definition_de}

\begin{bemerkung}
Die Kuratowski-Einbettung $\iota(x) = d(x,\cdot)$ ist stets eine isometrische
Einbettung in $L^\infty(M)$.  Die Definition ist unabhängig von der Wahl des
umgebenden Banach-Raums.
\end{bemerkung}

\begin{example}
\begin{enumerate}[label=(\roman*)]
  \item Für die runde $n$-Sphäre $S^n(R)$ vom Radius $R$:
        $\FillRad(S^n(R)) = R \arccos\!\bigl(-\tfrac{1}{n+1}\bigr)$
        (Katz \cite{Katz1983}).
  \item Für $\RP^n$ mit der runden Metrik vom Radius $R$:
        $\FillRad(\RP^n) = \tfrac{\pi R}{3}$.
\end{enumerate}
\end{example}

% =============================================================
\section{Wesentliche Mannigfaltigkeiten}
% =============================================================

\begin{definition_de}[Wesentliche Mannigfaltigkeit]
Eine geschlossene $n$-Mannigfaltigkeit $M$ heißt \emph{wesentlich}, falls
die klassifizierende Abbildung $f \colon M \to K(\pi_1(M), 1)$ die
Bedingung $f_*[M] \neq 0 \in H_n(K(\pi_1(M),1);\ZZ)$ erfüllt.
\end{definition_de}

\begin{example}
\begin{enumerate}[label=(\roman*)]
  \item Flächen vom Geschlecht $g \ge 1$ sind wesentlich.
  \item $\RP^n$ ist wesentlich.
  \item $S^n$ ($n \ge 2$) ist \emph{nicht} wesentlich.
  \item Asphärische Mannigfaltigkeiten (kontrahierbarer universeller Überlagerungsraum)
        sind wesentlich.
\end{enumerate}
\end{example}

% =============================================================
\section{Gromovs Systolische Ungleichung}
% =============================================================

\begin{satz}[Gromov, 1983 \cite{Gromov1983}]
\label{satz:gromov-systolisch}
Für jede wesentliche geschlossene Riemannsche $n$-Mannigfaltigkeit $(M,g)$ gilt
\[
  \sys_1(M,g)^n \le C_n \cdot \vol(M,g),
\]
wobei $C_n > 0$ nur von $n$ abhängt.
\end{satz}

Gromovs Beweis verläuft in zwei Hauptschritten:

\begin{satz}[Füllungsradius-Schranke \cite{Gromov1983}]
\label{satz:fuellrad-sys}
Für jede geschlossene Riemannsche $n$-Mannigfaltigkeit $(M,g)$ gilt
\[
  \FillRad(M,g) \le \frac{1}{2} \sys_1(M,g).
\]
\end{satz}

\begin{satz}[Volumen-Füllungs-Ungleichung \cite{Gromov1983}]
\label{satz:vol-fuellrad}
Für jede wesentliche geschlossene Riemannsche $n$-Mannigfaltigkeit $(M,g)$ gilt
\[
  \FillRad(M,g)^n \le C_n' \cdot \vol(M,g).
\]
\end{satz}

Satz~\ref{satz:gromov-systolisch} folgt unmittelbar durch Kombination von
Satz~\ref{satz:fuellrad-sys} und Satz~\ref{satz:vol-fuellrad}.

\begin{proof}[Beweisidee zu Satz~\ref{satz:fuellrad-sys}]
Sei $\gamma$ eine kürzeste nicht-kontrahierbare Schleife der Länge $L = \sys(M)$.
Falls $\FillRad(M) > L/2$, bleibt $\gamma$ in der $L/2$-Umgebung von
$\iota(M)$ in $L^\infty(M)$ nicht-kontrahierbar.  Ein Nerv-Argument
(Vietoris--Rips-Komplex bei Skala $\varepsilon \to L/2$) zeigt, dass die
Fundamentalklasse $[M]$ bereits bei Radius $\le L/2$ verschwindet --
ein Widerspruch. $\square$
\end{proof}

% =============================================================
\section{Pus Ungleichung für $\RP^2$}
% =============================================================

\begin{satz}[Pu, 1952 \cite{Pu1952}]
Für jede Riemannsche Metrik $g$ auf $\RP^2$ gilt
\[
  \sys(\RP^2, g)^2 \le \frac{\pi}{2} \vol(\RP^2, g),
\]
mit Gleichheit genau dann, wenn $(\RP^2, g)$ die runde Metrik konstanter
Krümmung ist (mit passender Normierung).
\end{satz}

\begin{proof}[Beweisidee]
Pus Beweis nutzt den Uniformisierungssatz: $g = e^{2u} g_{\mathrm{rund}}$.
Die Systole wird durch die konforme Länge einer nicht-kontrahierbaren Geodäten
nach unten beschränkt, und die Cauchy--Schwarz-Ungleichung liefert die Schranke.
Der Gleichheitsfall ergibt sich durch Analyse. $\square$
\end{proof}

Die Konstante $\pi/2$ heißt die \emph{Pu-Konstante}.

% =============================================================
\section{Loewners Torus-Ungleichung}
% =============================================================

\begin{satz}[Loewner, 1949, vgl.\ \cite{Pu1952}]
Für jede Riemannsche Metrik $g$ auf dem 2-Torus $\TT^2$ gilt
\[
  \sys(\TT^2, g)^2 \le \frac{2}{\sqrt{3}} \vol(\TT^2, g),
\]
mit Gleichheit genau dann, wenn $(\TT^2, g)$ (bis auf Skalierung) der flache
Torus $\RR^2/\Lambda$ für das hexagonale Gitter
$\Lambda = \ZZ(1,0) + \ZZ(\tfrac{1}{2},\tfrac{\sqrt{3}}{2})$ ist.
\end{satz}

\begin{proof}[Beweisidee]
Durch Uniformisierung: $g = e^{2u} g_\Lambda$ für eine flache Metrik $g_\Lambda$.
Die Ungleichung $\sys^2 \le C \vol$ folgt aus der Gitter-Geometrie:
Unter allen Rang-2-Gittern gegebenen Ko-Volumens wird der kürzeste Vektor
durch das hexagonale Gitter maximiert (Satz von Lagrange über
Gitterreduktion). $\square$
\end{proof}

% =============================================================
\section{Bergers Problem und scharfe Konstanten}
% =============================================================

\begin{definition_de}[Systolisches Verhältnis]
Das \emph{systolische Verhältnis} von $(M,g)$ ist
$\mathrm{SR}(M,g) = \sys(M,g)^n/\vol(M,g)$.
Das \emph{optimale systolische Verhältnis} ist $\mathrm{SR}(M) = \sup_g \mathrm{SR}(M,g)$.
\end{definition_de}

\begin{center}
\begin{tabular}{lll}
\toprule
Mannigfaltigkeit & $\mathrm{SR}(M)$ & Extremale Metrik \\
\midrule
$\TT^2$ & $2/\sqrt{3}$ & Hexagonaler Torus (Loewner) \\
$\RP^2$ & $\pi/2$ & Runde Metrik (Pu) \\
$\TT^3$ & ? & Offen (Berger) \\
$\RP^3$ & ? & Offen \\
\bottomrule
\end{tabular}
\end{center}

\begin{vermutung}[Bergers Problem über scharfe systolische Konstanten -- OFFEN]
Bestimme das optimale systolische Verhältnis $\mathrm{SR}(M)$ für alle
geschlossenen Riemannschen Mannigfaltigkeiten~$M$.  Insbesondere für
$\TT^n$ ($n \ge 3$) und $\RP^n$ ($n \ge 3$).
\end{vermutung}

% =============================================================
\section{Morse-Theorie für den Schleifenraum}
% =============================================================

Die Morse-Theorie des Energie-Funktionals
$E(\gamma) = \int_0^1 |\dot\gamma|^2 dt$ auf dem freien Schleifenraum
$\Lambda M$ ist ein mächtiges Werkzeug in der systolischen Geometrie.

\begin{satz}[Ljusternik--Schnirelmann für Geodäten]
Jede geschlossene Riemannsche Mannigfaltigkeit mit $\pi_1(M) \neq 1$
enthält eine \emph{geschlossene Geodäte} der Länge genau $\sys_1(M)$.
\end{satz}

\begin{satz}[Gromovs Taillen-Ungleichung \cite{Gromov2003waist}]
Sei $M$ eine geschlossene Riemannsche $n$-Mannigfaltigkeit und
$f \colon M \to \RR^k$ stetig.  Dann gibt es eine Faser $f^{-1}(t)$ mit
\[
  \sys(f^{-1}(t))^{n-k} \le C_{n,k} \cdot \vol_k(M)^{(n-k)/n}.
\]
\end{satz}

% =============================================================
\section{Sabourau--Rotman-Verbesserungen}
% =============================================================

\begin{satz}[Rotman, 2006 \cite{Rotman2006}]
Sei $M$ eine geschlossene Riemannsche $n$-Mannigfaltigkeit mit $\pi_1(M) \neq 1$.
Dann gilt für die kürzeste nicht-kontrahierbare geschlossene Geodäte $\gamma$:
\[
  \mathrm{Länge}(\gamma) = \sys(M) \le 2 \diam(M).
\]
\end{satz}

\begin{satz}[Sabourau, 2004 \cite{Sabourau2004}]
Für jede asphärische Fläche $(M^2, g)$ (d.\,h.\ Geschlecht $g \ge 1$) gilt
\[
  \sys(M)^2 \le \gamma(M) \cdot \vol(M),
\]
wobei $\gamma(M) > 0$ nur von der Topologie von $M$ abhängt.
\end{satz}

\begin{vermutung}[Optimales Füllungsverhältnis -- OFFEN]
Die scharfe Konstante in $\FillRad(M)^n \le C \cdot \vol(M)$ (für wesentliche $M$)
wird von einer spezifischen "runden" oder "regulären" Metrik realisiert.
Für $\RP^n$ ist der vermutete scharfe Wert $\tfrac{\pi}{3}$.
\end{vermutung}

% =============================================================
\section{Katz' Resultate zum Füllungsradius}
% =============================================================

\begin{satz}[Katz, 1983 \cite{Katz1983}]
Der Füllungsradius der runden $n$-Sphäre $S^n(R)$ vom Radius $R$ ist
\[
  \FillRad(S^n(R)) = R \arccos\!\left(-\frac{1}{n+1}\right).
\]
\end{satz}

\begin{satz}[Katz, 1983 \cite{Katz1983}]
Für $\RP^n$ mit der runden Metrik vom Radius $R$:
\[
  \FillRad(\RP^n) = \frac{\pi R}{3}.
\]
\end{satz}

\begin{korollar}
Für $\RP^2$ mit beliebiger Metrik $g$ mit $\vol(\RP^2, g) = \pi$ gilt:
\[
  \FillRad(\RP^2, g) \le \frac{\pi}{3},
\]
mit Gleichheit für die runde Metrik -- in Übereinstimmung mit Pus Ungleichung.
\end{korollar}

% =============================================================
\section{Gromovs Kompaktheitssatz}
% =============================================================

\begin{satz}[Gromov-Kompaktheit \cite{Gromov1981}]
Die Klasse der Riemannschen $n$-Mannigfaltigkeiten $(M,g)$ mit Ricci-Krümmung
$\mathrm{Ric}(M,g) \ge -(n-1)\kappa^2$ und Durchmesser $\diam(M,g) \le D$
ist präkompakt in der Gromov--Hausdorff-Topologie: Jede Folge besitzt eine
in diesem Sinn konvergente Teilfolge.
\end{satz}

Dieser Satz wird in der systolischen Geometrie verwendet, um Grenzwerte von
Folgen von Metriken zu untersuchen und das Verhalten von Systole und
Füllungsradius in degenerierenden Familien zu analysieren.

% =============================================================
\section{Offene Probleme in der Systolischen Geometrie}
% =============================================================

\begin{vermutung}[Scharfe Systole für $\TT^n$, $n \ge 3$]
Das optimale systolische Verhältnis des $n$-Torus $\TT^n$ ist
\[
  \mathrm{SR}(\TT^n) = \gamma_n^{n/2},
\]
wobei $\gamma_n$ die \emph{Hermite-Konstante} ist.
Diese ist für $n \le 8$ und $n = 24$ aus der Gittertheorie bekannt,
aber für allgemeines $n$ offen.
\end{vermutung}

\begin{vermutung}[Gromovs optimale Füllungsradius-Vermutung -- OFFEN]
Die scharfe Konstante in $\sys(M)^n \le C \cdot \vol(M)$ hängt mit der
Geometrie des regulären Simplex in $\RR^n$ zusammen.
Die genaue scharfe Konstante in Satz~\ref{satz:gromov-systolisch}
ist für $n \ge 3$ unbekannt.
\end{vermutung}

\begin{vermutung}[Systolische Freiheit -- OFFEN]
Für einfach zusammenhängende Mannigfaltigkeiten der Dimension $\ge 4$
kann das systolische Verhältnis $\mathrm{SR}(M,g)$ beliebig groß werden
(Gromovs \glqq systolische Freiheit\grqq{}).
Eine vollständige Charakterisierung der Mannigfaltigkeiten, die systolische
Freiheit aufweisen, steht aus.
\end{vermutung}

% =============================================================
\section{Fazit}
% =============================================================

Die systolische Geometrie, begründet von Loewner, Pu und Gromov, stellt
tiefe Verbindungen zwischen metrischer Geometrie, algebraischer Topologie
und der Theorie der Banach-Räume her.
Der Füllungsradius---das Maß dafür, wie tief eine Mannigfaltigkeit in ihre
Kuratowski-Vervollständigung eingebettet ist---erweist sich als
entscheidende Zwischengröße.

Die scharfen Konstanten in systolischen Ungleichungen sind nur für Flächen
($\TT^2$, $\RP^2$, Kleinsche Flasche) bekannt und in höheren Dimensionen
offen.  Das Zusammenspiel mit Gittertheorie, Hermite-Konstanten und der
Geometrie optimaler Kugelpackungen deutet auf tiefe Verbindungen mit der
Zahlentheorie und algebraischen Geometrie hin.

\begin{thebibliography}{99}

\bibitem{Gromov1983}
M.~Gromov,
\textit{Filling Riemannian manifolds},
J. Differential Geom. \textbf{18} (1983), Nr.~1, 1--147.

\bibitem{Pu1952}
P.~M. Pu,
\textit{Some inequalities in certain nonorientable Riemannian manifolds},
Pacific J. Math. \textbf{2} (1952), 55--71.

\bibitem{Katz1983}
M.~Katz,
\textit{The filling radius of two-point homogeneous spaces},
J. Differential Geom. \textbf{18} (1983), Nr.~3, 505--511.

\bibitem{Gromov1981}
M.~Gromov,
\textit{Structures métriques pour les variétés riemanniennes},
Textes Mathématiques \textbf{1}, CEDIC, Paris, 1981.

\bibitem{Gromov2003waist}
M.~Gromov,
\textit{Isoperimetry of waists and concentration of maps},
Geom. Funct. Anal. \textbf{13} (2003), 178--215.

\bibitem{Rotman2006}
R.~Rotman,
\textit{The length of a shortest geodesic net on a closed Riemannian manifold},
Topology \textbf{46} (2007), 343--356.

\bibitem{Sabourau2004}
S.~Sabourau,
\textit{Systoles des surfaces plates singulières de genre deux},
Math. Z. \textbf{247} (2004), 693--709.

\bibitem{BergerPanorama}
M.~Berger,
\textit{A Panoramic View of Riemannian Geometry},
Springer, Berlin, 2003.

\bibitem{KatzSabourau}
M.~Katz und S.~Sabourau,
\textit{Hyperelliptic surfaces are Loewner},
Proc. Amer. Math. Soc. \textbf{134} (2006), 1189--1195.

\bibitem{KatzBuch}
M.~G.~Katz,
\textit{Systolic Geometry and Topology},
Mathematical Surveys and Monographs \textbf{137},
American Mathematical Society, 2007.

\end{thebibliography}

\end{document}
